% !PRINTED 679
\weberpage{679}
%<<CONT
$= 0$, das eine $= 1$ annehmen, so ergiebt sich aus (3) irgend eine der $(\nu-1)$-reihigen Determinanten der Matrix
% [sic] The source cites (3), but the determinant specialized by the x_i choices is (4) on printed p. 678.
\[
\begin{matrix}
l_{1,1}, & l_{1,2}, & \ldots, & l_{1,\nu} \\
\cdots & \cdots & \cdots & \cdots \\
l_{\nu-1,1}, & l_{\nu-1,2}, & \ldots, & l_{\nu-1,\nu}
\end{matrix}
\]
und alle diese Determinanten haben also (vom Vorzeichen abgesehen) denselben Werth. Es ist daher gleichg\"ultig, welcher unter den $\nu$ conjugirten Werthen bei der Bildung des Regulators $L$ weggelassen wird.

Ist $\varepsilon_1, \varepsilon_2, \ldots, \varepsilon_{\nu-1}$ ein unabh\"angiges System von Einheiten, und sind $\lambda_1, \lambda_2, \ldots, \lambda_\nu$ die conjugirten Logarithmen irgend einer Einheit $\varepsilon$ in $\Omega$, so kann man die Zahlen $\xi_1, \xi_2, \ldots, \xi_{\nu-1}$ immer und nur auf eine Weise so bestimmen, dass
\weq{5}{
\begin{array}{l}
\xi_1 l_{1,1} + \xi_2 l_{2,1} + \cdots + \xi_{\nu-1} l_{\nu-1,1} = \lambda_1, \\
\xi_1 l_{1,2} + \xi_2 l_{2,2} + \cdots + \xi_{\nu-1} l_{\nu-1,2} = \lambda_2, \\
\cdots\cdots\cdots\cdots\cdots\cdots\cdots \\
\xi_1 l_{1,\nu} + \xi_2 l_{2,\nu} + \cdots + \xi_{\nu-1} l_{\nu-1,\nu} = \lambda_\nu
\end{array}
}
wird. Denn von diesen Gleichungen sind die ersten $\nu - 1$ ein System linearer Gleichungen mit nicht verschwindender Determinante, und die $\nu^{\mathrm{te}}$ Gleichung folgt aus den \"ubrigen, weil die Summe der conjugirten Logarithmen einer jeden Einheit verschwindet.

Das System der Zahlen $\xi_1, \xi_2, \ldots, \xi_{\nu-1}$ nennen wir das \gesperrt{Exponentensystem} der Einheit $\varepsilon$ in Bezug auf das System $\varepsilon_1, \varepsilon_2, \ldots, \varepsilon_{\nu-1}$, und es ist nun zu beweisen:

6. \gesperrt{Dass die Exponenten jeder Einheit $\varepsilon$ rationale Zahlen sind, deren Nenner eine gewisse endliche Grenze nicht \"ubersteigt, und die daher alle mit einem gemeinschaftlichen, nur von dem Systeme $\varepsilon_i$ abh\"angigen Nenner behaftet angenommen werden k\"onnen.}

Um dies einzusehen, hat man Folgendes zu erw\"agen:

1) Wenn wir die Gr\"ossen $\xi_1, \xi_2, \ldots, \xi_{\nu-1}$ wie Variable betrachten, und jede von ihnen, von den anderen unabh\"angig, von 0 bis 1 gehen lassen, so bleiben die linken Seiten von (5) in endlichen, nur von den Einheiten $\varepsilon_1, \varepsilon_2, \ldots, \varepsilon_{\nu-1}$ abh\"angigen Grenzen eingeschlossen. In diesen Grenzen m\"ussen also auch die conjugirten Logarithmen $\lambda$ der Einheit $\varepsilon$ bleiben, so lange
%CONT>>
